\SourceLeaf{294}%
\[
\begin{aligned}
t&=t^0(1-\tfrac{1}{6}\rho\rho rr-\tfrac{1}{3}\rho\rho t^0t^0)\\
\lambda&=\lambda^0(1-\tfrac{1}{6}\rho\rho s^0s^0-\tfrac{1}{3}\rho\rho t^0t^0)\\
\sigma&=\tfrac{1}{2}\rho vt^0(1-\tfrac{1}{12}\rho\rho rr-\tfrac{1}{4}\rho\rho s^0s^0-\tfrac{1}{4}\rho\rho t^0t^0)\\
\tau&=\tfrac{1}{2}\rho vs^0(1+\tfrac{5}{12}\rho\rho rr-\tfrac{1}{2}\rho\rho s^0s^0)
\end{aligned}
\]
für \(t\) und \(\lambda\) auf die fünfte, für \(\sigma\) und \(\tau\) auf die sechste Ordnung (ausschl.) genau. Noch bequemer und eben so genau ist es, hiebei sogleich die Logarithmen zu gebrauchen, wodurch die Formeln, wenn man zur Abkürzung das Product der Grösse \(\tfrac{1}{12}\rho\rho\) in den Modulus der briggischen Logarithmen mit \(\mu\) bezeichnet, folgende Gestalt erhalten:
\[
\begin{aligned}
\log s&=\log s^0+4\mu rr-4\mu s^0s^0\\
\log t&=\log t^0-2\mu rr-4\mu t^0t^0\\
\log\lambda&=\log\lambda^0-2\mu s^0s^0-4\mu t^0t^0\\
\log\sigma&=\log\tfrac{1}{2}\rho vt^0-\mu rr-3\mu s^0s^0-3\mu t^0t^0\\
\log\tau&=\log\tfrac{1}{2}\rho vs^0+5\mu rr-6\mu s^0s^0
\end{aligned}
\]
\indent Diese fünf Formeln in Verbindung mit den vorhergehenden für \(s^0,t^0,\lambda^0\) bilden eine fünfte Auflösungsart, deren eigenthümliches es ist, dass genäherte Werthe der Grössen \(s,t,\lambda,\sigma,\tau\) durch kleine sehr leicht zu berechnende an den Logarithmen anzubringende Correctionen zu scharfen erhoben werden. Die hiebei vorkommenden constanten Logarithmen sind
\[
\begin{aligned}
\log\rho&=4{,}6855748668&&(-10)\\
\log\tfrac{1}{2}\rho&=4{,}3845448712&&(-10)\\
\log\mu&=7{,}9297527989&&(-20)
\end{aligned}
\]
oder wenn jene Correctionen sofort als Einheiten der siebenten Decimale erscheinen sollen
\[
\log\mu=4{,}9297527989\qquad(-10)
\]
von welchen Logarithmen jedoch hier nur die ersten Ziffern zur Anwendung kommen.
\Article{17}
Viel einfacher lassen sich aber die Relationen zwischen den Grössen \(r,S,S',T,T',\lambda\) ausdrücken, wenn man von dem Mittel der beiden Breiten
